\documentclass[11pt, oneside]{article}
\usepackage[margin=1in]{geometry}
\usepackage{booktabs}
\usepackage{hyperref}
\usepackage{enumitem}
\usepackage{float}
\usepackage{caption}
\usepackage{subcaption}
\usepackage[numbers]{natbib}
\usepackage{graphicx}
\usepackage{amssymb}
\usepackage{array}
\usepackage{multirow}

\title{CS780 Capstone Project:\\
OBELIX: the Warehouse Robot}

\author{
    FirstName $<$Middle Name if any$>$ LastName  \\
    roll no.\\
   {\tt \{name\}@iitk.ac.in}\\
{Indian Institute of Technology Kanpur (IIT Kanpur)}
}

\date{}

\begin{document}

\maketitle

\begin{abstract}
We tackle OBELIX, a warehouse robot POMDP task where an agent must find, attach to, push, and unwedge a grey box to the arena boundary using only binary sonar and IR sensors.
We explored Double DQN, vanilla PPO, PPO with frame stacking, PPO with LSTM, a three-policy hierarchical architecture, GRU-based multi-head PPO, and a final Explicit Phase Belief (EPB) agent that replaces recurrent memory with a hand-crafted non-decaying 16-dimensional belief state.
A key finding was that reward shaping and encoded observations hurt performance compared to raw reward PPO.
The best codabench score of $+157.4$ was achieved by the GRU-based agent trained without reward shaping.
\end{abstract}


\section{Competition Result}
\textbf{Codabench Username:} \_\_\_\_\_\_\_\_\_\_\_ \\
\textbf{Final leaderboard rank (Test Phase):} \_\_\_\_\_\_\_\_\_\_\_ \\
\textbf{Leaderboard rank -- Level 1 (static):} \_\_\_\_\_\_\_\_\_\_\_ \\
\textbf{Leaderboard rank -- Level 2 (blinking):} \_\_\_\_\_\_\_\_\_\_\_ \\
\textbf{Leaderboard rank -- Level 3 (moving + blinking):} \_\_\_\_\_\_\_\_\_\_\_ \\

\vspace{0.5em}
\textbf{Number of valid submissions (Development and Test):}

\begin{table}[H]
\centering
\begin{tabular}{lcccc}
\toprule
Phase & Agent & Mean Score & Std & Runs \\
\midrule
Dev (Feb 11) & agent\_template & $-13931.0$ & $8342.4$ & 3 \\
Dev (Feb 11) & agent\_template & $-16100.7$ & $10503.0$ & 3 \\
Dev (Mar 22) & agent (PPO, raw obs) & $+229.9$ & $1506.3$ & 10 \\
Dev (Mar 22) & agent (PPO, raw obs) & $+539.9$ & $1539.9$ & 10 \\
Dev (Mar 22) & agent (PPO, raw obs) & $+537.2$ & $1537.2$ & 10 \\
Dev (Mar 22) & agent (PPO, shaped) & $-111868.2$ & $93154.5$ & 10 \\
Dev (Mar 22) & agent (PPO, shaped+walls) & $-136026.6$ & $90449.6$ & 10 \\
Dev (Mar 29) & agent (Phase 2 final) & $-238.6$ & $2300.1$ & 10 \\
Dev (Apr 5) & agent\_three\_policy & $-31641.8$ & $67989.9$ & 10 \\
Dev (Apr 7) & agent\_three\_policy & $-73666.1$ & $92732.4$ & 10 \\
Dev (Apr 7) & agent\_three\_policy (walls) & $-158772.2$ & $73982.4$ & 10 \\
Dev (Apr 7) & agent\_three\_policy\_latent & $-114049.1$ & $94840.9$ & 10 \\
Dev (Apr 10) & agent\_ppo\_gru\_un & $-394.2$ & $4092.0$ & 10 \\
Dev (Apr 10) & agent\_ppo\_gru\_un (walls) & $+157.4$ & $1398.6$ & 10 \\
\bottomrule
\end{tabular}
\caption{All development submissions. Scores reported by codabench (scaling factor $\times 5$, 1000 max steps).}
\label{tab:submissions}
\end{table}


\section{Introduction and Problem Description}

OBELIX is a simulated warehouse robot task that requires an autonomous agent to navigate a 2D arena, locate a grey box, attach to it by approaching head-on, push the box to any boundary wall, and handle cases where the box becomes wedged in a corner (the ``unwedge'' phase).
The challenge is designed to be a realistic Partially Observable Markov Decision Process (POMDP): the robot has \emph{no global position or orientation}, and can only observe 18 binary sensor bits: 4 left sonar readings (near/far $\times$ 2 beams), 8 front sonar readings (near/far $\times$ 4 beams), 4 right sonar readings (near/far $\times$ 2 beams), one IR proximity sensor (object at $<4$~inches), and one stuck flag (motor current exceeded threshold).

The task unfolds in three difficulty levels:
\begin{itemize}[noitemsep]
    \item \textbf{Level 1 (static):} The box remains stationary. The robot must find and push it to a wall.
    \item \textbf{Level 2 (blinking):} The box periodically disappears and reappears. Sensors may return false negatives, breaking naive reactive strategies.
    \item \textbf{Level 3 (moving + blinking):} The box drifts autonomously in addition to blinking, making prediction and interception necessary.
\end{itemize}
An optional wall-obstacle variant places internal walls in the arena, requiring collision-aware navigation.

The POMDP nature makes this problem hard: the box may be invisible for hundreds of consecutive steps, the robot cannot localise itself, the phase transitions (find $\to$ push $\to$ unwedge) must be inferred purely from sensor evidence, and exploration is structureless without any global map.
The objective is to maximise mean cumulative reward (per episode, scaled by 5) across all difficulties, with a maximum of 1000 steps per episode.


\section{Approach and Model Evolution}

\subsection{Phase 1: Double DQN Baseline}

The starting point was a Double DQN (DDQN) agent trained on the raw 18-dimensional binary observation.
The network was a three-layer MLP (18~$\to$~64~$\to$~64~$\to$~5) with $\epsilon$-greedy exploration.
A vectorised environment (\texttt{vec\_env.py}) was written to run $N$ OBELIX instances in parallel subprocesses for faster data collection.
A \texttt{RewardShaper} (\texttt{reward\_shaper.py}) was also implemented, adding sensor-evidence bonuses during the find phase, a forward-progress bonus during push, an escalating stuck penalty, and a recovery bonus.
A \texttt{BeliefStateEncoder} (\texttt{state\_encoder.py}) was implemented to enrich the 18-bit observation with sector strengths, direction summaries, temporal memory flags, and transition cues, producing a 49-dimensional per-frame core which could be further stacked across $k$ timesteps.

DDQN with shaped rewards and encoded observations produced poor results on local evaluation; the results from the early template agent on codabench confirmed this ($-13931$, $-16100$, Table~\ref{tab:submissions}).

\subsection{Phase 2: PPO and the Surprising Benefit of Simplicity}

Switching to Proximal Policy Optimisation (PPO) with a shared-trunk MLP actor-critic (two Tanh hidden layers, 128 units each) and raw 18-dimensional observations showed a dramatic improvement.
When trained \emph{without} reward shaping and \emph{without} encoded observations, codabench mean scores reached $+229.9$ and then $+539.9$.
By contrast, enabling the \texttt{RewardShaper} and \texttt{BeliefStateEncoder} on the same PPO architecture led to scores of $-111{,}868$ and $-136{,}027$ -- roughly three orders of magnitude worse.

This result was unexpected and informative: the hand-crafted shaping terms and engineered features were actively misleading the agent, either introducing conflicting gradients or preventing the policy from learning the actual task reward signal.
The conclusion drawn was that for a sparse-reward POMDP of this type, \emph{clean raw observations and unmodified environment rewards are preferable} to dense shaping.

PPO with frame stacking (stacking $k$ raw observations) was also evaluated.
It gave modest improvements on local evaluation but did not clearly outperform the single-frame baseline on codabench.

An LSTM-augmented policy (\texttt{train\_ppo\_lstm.py}) was tried next as a natural extension for handling partial observability.
The LSTM maintained a hidden state across timesteps and could, in principle, remember the last box direction across blank sensor steps.
However, in practice the LSTM agent also failed to match the raw-PPO baseline, particularly when combined with shaped rewards.
By the end of Phase 2, the best performing checkpoint remained the plain PPO agent with raw observations and unshaped rewards, scoring $-238.6$ in the final Phase~2 submission.

\subsection{Phase 3: Three-Policy Architecture and GRU Recovery}

Motivated by reading the OBELIX paper~\cite{obelix}, Phase 3 explored a behaviour-decomposed architecture matching the paper's three-behaviour hierarchy:
\[
\text{unwedge} \succ \text{push} \succ \text{find}
\]
A \texttt{BehaviorManager} (\texttt{behavior\_manager.py}) selected the active sub-policy at each step based on sensor evidence.
Three independent PPO actor-critic heads were trained, each receiving the same 18-bit (or encoded) observation.
An \texttt{AntiSpinPenaltyTracker} (\texttt{anti\_spin\_penalty.py}) was added to penalise repetitive rotation during the find phase.

Despite alignment with the paper's intuition, the three-policy agent performed poorly in practice.
Codabench scores ranged from $-31{,}641$ to $-158{,}772$ across variants.
Switching to a latent-space version (\texttt{agent\_three\_policy\_latent}) did not help, yielding $-114{,}049$.
The likely cause was policy interference: the sharp handoffs between sub-policies produced inconsistent reward signals, and the naive FSM phase detector frequently mis-classified the current phase.

Recognising the failure of hard phase-switching, the next attempt (\texttt{ppo\_gru.py}, \texttt{agent\_ppo\_gru\_un.py}) replaced the three separate heads with a unified agent containing:
\begin{itemize}[noitemsep]
    \item An MLP actor-critic for the find and push phases (shared backbone, 18 $\to$ 128 $\to$ 128, Tanh).
    \item A separate GRU-based actor-critic for the unwedge phase, where temporal context matters most for detecting that the box is wedged.
    \item A rule-based phase selector that latches on IR attachment.
\end{itemize}
This agent was trained without any reward shaping, directly on the raw environment reward.
It achieved codabench scores of $-394.2$ and $+157.4$, the best result of the project.

\subsection{EPB: Explicit Phase Belief Agent (Final Approach)}

The GRU agent suffered from a known failure mode: after many consecutive blank sensor steps (common in Level~2 and Level~3), the GRU hidden state is washed toward zero by the recurrent dynamics, erasing the last-known box direction.
Phase mis-detection by the FSM also caused spurious transitions.

To address both issues, the \textbf{Explicit Phase Belief (EPB)} agent (\texttt{train\_ppo\_epb.py}, \texttt{agent\_epb.py}) was designed around a hand-crafted 16-dimensional belief state (\texttt{compact\_belief.py}):

\begin{itemize}[noitemsep]
    \item \textbf{Non-decaying last-visible direction:} Fields \texttt{last\_visible\_left/front/right/ir} update \emph{only when a sensor fires} and never decay on blank steps.
    \item \textbf{Staleness counter:} \texttt{steps\_since\_visible} counts up (capped at 500); the agent knows how stale its last sighting is.
    \item \textbf{Attachment latch:} \texttt{attached} latches True on two consecutive IR activations and is never reset within an episode.
    \item \textbf{Stuck severity, recovery edge, exploration direction:} additional scalars for nuanced behaviour.
\end{itemize}

The 39-dimensional input to the policy is:
\[
[\underbrace{\text{raw obs}}_{18} \;\|\; \underbrace{\text{belief vector}}_{16} \;\|\; \underbrace{\text{FSM one-hot}}_{5}]
\]
The FSM one-hot encodes a reactive hand-coded action suggestion (orient toward nearest sensor, spin if blind, push forward if attached) as a soft prior.

The network is a \textbf{dual-head actor-critic}: a shared MLP trunk (39~$\to$~128~$\to$~128, Tanh) feeds both a \emph{find} actor head and a \emph{push} actor head, plus a learned sigmoid gate supervised by \texttt{env.enable\_push} and a shared critic head.
The combined logits are $\mathbf{l} = (1 - g)\,\mathbf{l}_\text{find} + g\,\mathbf{l}_\text{push}$ where $g$ is the gate output.

Training used a 5-stage arena curriculum to ensure dense push-head gradients early:
\begin{enumerate}[noitemsep]
    \item 300~px arena, static box, 800 episodes
    \item 400~px arena, static box, 1000 episodes
    \item 500~px arena, static box, 1200 episodes
    \item 500~px arena, static + walls, 1200 episodes
    \item 500~px arena, blinking + walls, 1000 episodes
    \item 500~px arena, moving + blinking + walls, 800 episodes
\end{enumerate}


\section{Experiments and Implementation Details}

All training was done on a GPU-accelerated machine (CUDA).
The PPO hyperparameters used across most variants are listed in Table~\ref{tab:hparams}.

\begin{table}[H]
\centering
\begin{tabular}{lcc}
\toprule
Hyperparameter & PPO (Phase~2) & EPB (Phase~3) \\
\midrule
Max steps/episode & 1000 & 600--1000 (curriculum) \\
Rollout length & 512 & 512 \\
Parallel environments & 8 & 8 \\
Discount $\gamma$ & 0.99 & 0.99 \\
GAE $\lambda$ & 0.95 & 0.95 \\
PPO clip $\epsilon$ & 0.2 & 0.2 \\
Entropy coefficient & 0.01 & 0.01 \\
Value loss coefficient & 0.5 & 0.5 \\
Learning rate & $3 \times 10^{-4}$ & $3 \times 10^{-4}$ \\
Optimiser & Adam & Adam \\
PPO epochs per update & 4 & 4 \\
Mini-batch size & 256 & 256 \\
Hidden layer size & 128 & 128 \\
Observation dimension & 18 & 39 (belief-augmented) \\
Gate auxiliary loss weight & -- & 0.5 \\
Total episodes & 3000 & 6000 \\
\bottomrule
\end{tabular}
\caption{Hyperparameters for the two main approaches.}
\label{tab:hparams}
\end{table}

For the GRU-based agent (\texttt{ppo\_gru.py}), the find and push policies used the MLP backbone above; the unwedge policy used a GRU with hidden size 64, unrolled over the episode.
The action set was $\{\text{L45}, \text{L22}, \text{FW}, \text{R22}, \text{R45}\}$ for all policies.

Local evaluation used 10 independent seeds per checkpoint, matching the codabench protocol.
Learning curves were monitored to detect overfitting to the curriculum stage and to decide when to advance to the next stage.


\section{Results}

\begin{table}[H]
\centering
\begin{tabular}{lcc}
\toprule
Model & Codabench Mean Score & Notes \\
\midrule
DDQN (template agent) & $-13{,}931$ & Phase~1 baseline, 3 runs, 200 steps \\
PPO, raw obs (best) & $+539.9$ & Phase~2, no shaping, 10 runs \\
PPO, shaped + encoded & $-111{,}868$ & Phase~2, shaping hurts badly \\
PPO, Phase~2 final & $-238.6$ & Phase~2 submission \\
Three-policy PPO & $-31{,}641$ & Phase~3, paper-inspired hierarchy \\
Three-policy PPO (walls) & $-158{,}772$ & Phase~3, worst result \\
Three-policy latent & $-114{,}049$ & Phase~3 variant \\
GRU multi-head PPO & $+157.4$ & Phase~3 \textbf{best overall}, with walls \\
\bottomrule
\end{tabular}
\caption{Summary of model performance. Scores are codabench mean reward (scaled $\times 5$, 1000 max steps, 10 runs unless noted).}
\label{tab:results}
\end{table}

The GRU multi-head agent (\texttt{agent\_ppo\_gru\_un}) achieved the best codabench score of $+157.4$ on the walls variant of the development set.
This corresponds to the agent successfully completing the full find-push-unwedge sequence in some fraction of runs.

The most striking result of the project was the \emph{inverted performance} when reward shaping was introduced: the scored dropped from $+539.9$ to $-111{,}868$ on the exact same PPO architecture.
This suggests that the shaped reward terms (sensor bonuses, stuck penalties) created local optima that trapped the agent away from the sparse terminal reward.


\section{Error Analysis and Discussion}

\paragraph{Level 1 (static box).}
Even on the simplest level, the agent sometimes performs a prolonged spin without navigating toward the box, indicating that the find policy's exploration strategy is sub-optimal when the box is far from the starting position.
The GRU agent mitigates this by maintaining directional memory, but can still lose the box direction after extended blank-sensor periods.

\paragraph{Level 2 (blinking box).}
When the box blinks (disappears for several steps), agents that rely on reactive sensors immediately lose track.
The GRU's hidden state decays toward zero during blank steps -- the ``memory washout'' problem -- causing the agent to behave as if it has never seen the box.
The EPB agent was designed to fix this by storing the last-visible direction in persistent fields, but due to time constraints this agent was not submitted to the final test phase.

\paragraph{Level 3 (moving + blinking).}
A moving box requires the agent to predict and intercept rather than follow.
With only binary sonar readings and no velocity information, this is extremely difficult.
The GRU agent partially compensates via temporal pattern matching in the hidden state, but fails when the box drifts far enough to escape all sensor cones between sightings.

\paragraph{With walls.}
Internal wall obstacles cause the robot to get stuck more frequently.
The \texttt{STUCK} bit (bit~17) provides evidence of collision, but recovering from corner wedging requires backward-and-turn manoeuvres that are not explicitly rewarded.
Interestingly, the GRU agent scored \emph{higher} on the walls variant ($+157.4$) than without walls ($-394.2$), suggesting that walls may sometimes help corral the box toward a boundary.

\paragraph{Limitations and future work.}
\begin{itemize}[noitemsep]
    \item The EPB agent with the full 6-stage curriculum was not fully trained and evaluated on codabench before the deadline.
    \item A recurrent state-space model (RSSM, \texttt{train\_ppo\_rssm.py}) was begun but not completed.
    \item Hyperparameter search was largely manual; Bayesian optimisation could improve results significantly.
    \item The three-policy architecture, inspired directly by the OBELIX paper, was surprisingly ineffective. A softer mixture-of-experts gating might be preferable to hard FSM switching.
\end{itemize}


\section{Conclusion}

The main takeaway of this project is counter-intuitive: in a sparse-reward POMDP, \emph{less engineering is more}.
Hand-crafted reward shaping and feature encoding consistently hurt performance, while raw PPO trained directly on the environment reward achieved scores three orders of magnitude better.
This likely reflects reward hacking: the agent optimises the artificial shaped signal at the expense of the true objective.

The best agent (GRU multi-head PPO, $+157.4$) demonstrates that memory is beneficial for this POMDP -- but the type of memory matters.
Explicit, non-decaying belief states (EPB) are a principled alternative to GRU washout, and the architecture design appears sound; completing the training run is the primary remaining step.

Future directions include: training with model-based RL (RSSM / Dreamer) to learn a latent world model from binary observations, applying curriculum learning more aggressively, and exploring soft gating (mixture of experts) rather than hard FSM phase detection.


\section*{LLM Usage Declaration}
\begin{itemize}
\item \textbf{Which Large Language Model (LLM) did you use?} \\
Claude (Anthropic) via Claude Code CLI.

\item \textbf{In which part did you use the LLM, and how?}

\begin{tabular}{l|l|p{6cm}}
\toprule
Part & Used LLM? & Purpose of Using LLM \\
\midrule
Understanding concepts & Yes & Clarifying PPO/POMDP concepts and GRU memory dynamics \\
Ideation / brainstorming & Yes & Discussing EPB belief state design and curriculum structure \\
Debugging small code snippets & Yes & Fixing shape mismatches and vectorised-env bugs \\
Plotting / visualization code & No &   \\
Grammar / writing style & Yes & Improving report prose and structure \\
Report structure / section ideas & Yes & Drafting section outlines \\
Other & Yes & Generating training scripts (\texttt{train\_ppo\_epb.py}, \texttt{compact\_belief.py}) \\
\bottomrule
\end{tabular}

\item \textbf{Other sources consulted:} \\
OBELIX competition paper~\cite{obelix}; Schulman et al.\ PPO paper~\cite{ppo}; Stable-Baselines3 documentation for PPO implementation reference.

\item \textbf{Personal contribution:} \\
The core experimental journey -- selecting which algorithms to try, recognising that shaped rewards hurt performance, designing the three-policy vs.\ single-agent ablation, devising the non-decaying EPB belief state, and running all training and evaluation -- was driven by the student.
The LLM assisted with implementation and writing; all architectural decisions, observations, and conclusions are the student's own.
\end{itemize}

\vspace{2cm}

\bibliographystyle{ieeetr}
\bibliography{references}

\end{document}
